\chapter{Binäre Zerlegung von $K$ ist verlustfrei}

\bigskip
\begin{lemma}\label{lem:K_bin_ist_verlustfrei}
    Die entitätsbezogene Relation
    \[
        K(e, k_i, k_j), \ i \neq j
    \]
    lässt sich verlustfrei in die entitätsbezogenen Relationen
    \[
        R_{i}(e, k_i),\quad R_{j}(e, k_j)
    \]
    zerlegen.
\end{lemma}

\medskip
\begin{proof}
    Sei
    \[
        R_{i} = \pi_{(e,k_i)}(K),
        \qquad
        R_{j} = \pi_{(e,k_j)}(K).
    \]

    \noindent
    Wir zeigen, dass
    \[
        \njoin\bigl(\pi_{R_{i}}(r), \pi_{R_{j}}(r)\bigr)
        =
        \pi_{K}(r).
    \]

    \bigskip
    \noindent
    Da es sich um eine binäre Zerlegung handelt, kann der Beweis anhand der funktionalen Eigenschaften erfolgen~\cite[S. 343]{EN11}:
    \[
        \begin{alignedat}{3}
            ((R_{i} \cap R_{j}) &\rightarrow (R_{i} - R_{j}))
            = (\{e\} \rightarrow \{k_i\}) \in F_{K}^+\\
            ((R_{i} \cap R_{j}) &\rightarrow (R_{j} - R_{i}))
            = (\{e\} \rightarrow \{k_j\}) \in F_{K}^+
        \end{alignedat}
    \]
\end{proof}